\songtitle{Confusión}

% -> stories.tex (Children's songs and small creatures); original.tex (Other Narratives)
\tag*{2026}\tag{folk-pop}\tag{original-lyrics}\tag{spanish-lyrics}
\begin{notes}
Written by Carlos Thompson -- a barnyard in confusion where each stanza names the same animal by its regional Spanish names: the capybara (\emph{chigüire}, \emph{chigüiro}, \emph{carpincho}), the guinea pig (\emph{curí}, \emph{cuy}, \emph{cobaya}), the turkey (\emph{guajalote}, \emph{pisco}, \emph{guanajo}) and the opossum (\emph{chucha}, \emph{fara}, \emph{tlacuache}). Retrieved from Dvigitt's Suno page (V6-MINI, 3 October 2026).
\end{notes}
\begin{style}
Folk-pop with a vintage 70s mix, lo-fi room ambience, tape saturation, plate reverb, and a low-pass dubstep sub-bass drop; acoustic guitar strum, clapping percussion, accordion stabs, and children’s chorus; bouncing 6/8 waltz; male group vocals trade Spanish narrative lines and playful English call-and-response, building to a festive singalong hook.
\end{style}
\begin{lyrics}
\songpart{Verse 1}
El chigüire se levanta\\
el chigüiro pone la mesa\\
el carpincho pasa el mate\\
y el capibara pregunta: what?

\songpart{Chorus}
confusión en la granja\\
confusión en el campo\\
y todos responden\\
lo que todos callan

\songpart{Verse 2}
un curí pasa corriendo\\
el cuy dice que no pasa nada\\
la cobaya mira con recelo\\
y el conejillo de indias sigue perdido

\songpart{Chorus}
confusión en la granja\\
confusión en el campo\\
y todos responden\\
lo que todos callan

\songpart{Verse 3}
el guajalote se para altivo\\
el pisco lo mira de reojo\\
el guanajo se desentiende\\
y el pavo los ve y se pierde

\songpart{Chorus}
confusión en la granja\\
confusión en el campo\\
y todos responden\\
lo que todos callan

\songpart{Verse 4}
la chucha baja del árbol\\
una fara pasa de largo\\
el tlacuache se ríe de eso\\
y la zarigüeya se hace la tonta

\songpart{Final Chorus}
confusión en la granja\\
confusión en el campo\\
confusión en el monte\\
confusión en el prado
\end{lyrics}

